\documentclass[12pt]{article}
\usepackage[T1]{fontenc}
\usepackage[utf8]{inputenc}
\usepackage{lmodern}
\usepackage{textcomp}
\usepackage{enumitem}
\usepackage{geometry}
\geometry{margin=1in}
\begin{document}
\begin{center}
Answer Key - Business Administration - Graduate Level Assessment 23 \
\small (Answers are ordered exactly as in the question paper.)
\end{center}

\section*{Multiple Choice}
\begin{enumerate}[label=\arabic*.]
\item \textbf{Answer:} E
\\ \textit{Options:} A) \$45.67; B) \$35.28; C) \$53.60; D) \$48.99; E) \$41.52
\\ \textit{Note:} The correct answer of $41.52 is derived by calculating the interest on the loan using the formula I = PRT, where the principal (P) is $1,262.77, the rate (R) is 0.08 annually, and the time (T) is the exact number of days from March 15 to August 12, which is 120 days or 120/365 years, resulting in an interest amount of \$41.52.
\item \textbf{Answer:} E
\\ \textit{Options:} A) Hold objectives.; B) Expansion objectives.; C) Market objectives.; D) Competitive objectives.; E) Niche objectives
\\ \textit{Note:} Niche objectives are correct because they allow firms with limited financial resources to focus on a specific segment of the market, enabling them to compete effectively against a dominant competitor by catering to specialized needs rather than attempting to compete broadly.
\item \textbf{Answer:} A
\\ \textit{Options:} A) \$12.00; B) \$14.72; C) \$6.44; D) \$16.64; E) \$32
\\ \textit{Note:} The correct answer is $12.00 because Mr. Sarbo's policy was canceled after approximately 2.5 months of coverage, entitling him to a prorated refund for the unused portion of the annual premium, which results in a refund of $12 after accounting for the premium paid for the months of coverage used.
\item \textbf{Answer:} D
\\ \textit{Options:} A) \$190.20; B) \$152.35; C) \$198.75; D) \$174.23; E) \$210.00
\\ \textit{Note:} The correct answer is $174.23 because the selling price of $226.50 represents a 130\% of the cost, so dividing $226.50 by 1.30 yields the cost, which is $174.23.
\item \textbf{Answer:} D
\\ \textit{Options:} A) \$256.97; B) \$275.00; C) \$300.00; D) \$262.11; E) \$258.00
\\ \textit{Note:} The correct answer of $262.11 is derived by first calculating the base tax as $6,640 divided by $100 multiplied by the tax rate of $3.87, resulting in $256.99, and then adding a 2% collector's fee of $5.12 to that amount, leading to the total of \$262.11.
\end{enumerate}

\section*{True / False}
\begin{enumerate}[label=\arabic*.]
\setcounter{enumi}{5}
\item \textbf{Answer:} True
\\ \textit{Options:} A) True; B) False
\\ \textit{Note:} The statement is true because the final amount of \$338.27 accurately reflects the total after all applicable discounts have been deducted from the original invoice amount.
\item \textbf{Answer:} True
\\ \textit{Options:} A) True; B) False
\\ \textit{Note:} The net price of the Speedway Racing Set after applying a 10\% discount followed by a 20\% discount on the list price of $32 is calculated to be $29.12, confirming the accuracy of the statement.
\item \textbf{Answer:} True
\\ \textit{Options:} A) True; B) False
\\ \textit{Note:} The combined effect of a 20\% discount followed by a 10\% discount on the reduced price results in a total discount of approximately 28\% from the original price due to the sequential application of percentage discounts on a diminishing base.
\item \textbf{Answer:} False
\\ \textit{Options:} A) True; B) False
\\ \textit{Note:} Cognitive dissonance refers to the psychological discomfort experienced when a person holds conflicting beliefs or attitudes, rather than a process of evaluating products in relation to their self-concept.
\item \textbf{Answer:} False
\\ \textit{Options:} A) True; B) False
\\ \textit{Note:} The statement is false because a 30\% markup on selling price means the cost is actually 70\% of the selling price, so the cost would be $200,000 / 1.30, which is approximately $153,846, not \$180,000.
\end{enumerate}

\section*{Long Form Response}
\begin{enumerate}[label=\arabic*.]
\setcounter{enumi}{10}
\item \textbf{Answer:} George Mason's financing terms for the grand piano purchase reveal an interest rate of 12.7\% on the installment payments. This rate suggests a relatively high cost of borrowing, which can significantly impact the total amount paid over the life of the loan. In a broader financial context, such an interest rate may indicate higher risk associated with the financing arrangement, potentially reflecting the borrower's creditworthiness or the lender's risk assessment. Additionally, a 12.7\% interest rate may discourage consumers from financing larger purchases, prompting a reevaluation of personal budgeting and financial planning strategies in order to mitigate the effects of high-interest debt. Overall, understanding the implications of this interest rate is crucial for making informed financial decisions.
\\ \textit{Note:} correct because it accurately identifies the high interest rate of 12.7\% as indicative of the cost of borrowing and potential risk factors involved in the financing terms, while also linking these aspects to broader implications for creditworthiness and lender risk assessment.
\item \textbf{Answer:} Earnings per share (EPS) for Velco Corporation can be calculated by dividing the total profit by the number of shares outstanding. In this case, with a profit of $5,250 and 1,250 shares, the EPS is $5,250 / 1,250 = \$6.00. This figure is significant for shareholders as it indicates the company's profitability on a per-share basis, providing a clear measure of financial performance. A higher EPS often suggests better profitability, which can enhance shareholder confidence and potentially lead to an increase in stock price. Additionally, EPS is a critical metric used by investors to compare the financial health of Velco Corporation relative to its competitors in the market.
\\ \textit{Note:} correct because it accurately calculates the earnings per share (EPS) by dividing the total profit by the number of shares outstanding, demonstrating how EPS serves as a key indicator of the company's profitability and its implications for shareholder confidence and investment potential.
\end{enumerate}

\vfill
\noindent\textit{End of Answer Key}
\end{document}